Programming language interoperability is a crucial aspect of software development. A common challenge arises when bridging languages with considerable semantic differences. 
With the increasing expressiveness of modern programming languages, interoperability systems should facilitate more expressive communication. In this paper, we survey the capabilities of existing FFI-based interoperability systems and assess the viability of more expressive interoperability by extending an existing tool to support bridging polymorphic types. Through our tool-design survey and practical case study, we demonstrate that creating and applying more expressive interoperability between languages exhibiting significant semantic gaps is feasible. However, it requires a nuanced understanding of the languages involved and presents key trade-offs that must be considered.